% =========================
% report/evaluation.tex  (DROP-IN REPLACEMENT; VAE + cWGAN-GP)
% =========================

\subsection{Evaluation Objectives}
Evaluation is designed to answer three scientific questions under controlled, reproducible conditions:

\begin{enumerate}
  \item \textbf{VAE capacity and inference quality:}
  Does the modified VAE (e.g., flexible prior such as VampPrior, hierarchical latent design when enabled, and/or architectural capacity changes) improve latent utilization, sample quality, and likelihood-related metrics relative to a baseline VAE with a simple prior?

  \item \textbf{Posterior collapse behavior:}
  How do KL dynamics, active units, and aggregate posterior statistics differ between baseline and proposed models, and do these differences indicate reduced collapse risk or more effective latent usage?

  \item \textbf{Augmentation effectiveness:}
  Under a limited labeled-data regime, does conditional WGAN-GP augmentation improve downstream classifier generalization \emph{on real test data} without contaminating the classifier training split or introducing systematic class confusion?
\end{enumerate}

\subsection{General Evaluation Protocol and Reproducibility Controls}
\paragraph{Fixed splits and no leakage.}
For the augmentation experiment, all subset indices defining $\mathcal{D}_{\mathrm{GAN}}$ and $\mathcal{D}_{\mathrm{CLS}}$ are stored to disk and reused across runs. We enforce:
\[
\mathcal{D}_{\mathrm{GAN}} \cap \mathcal{D}_{\mathrm{CLS}} = \emptyset,
\]
and log the intersection size (must be zero) as an integrity check. For all experiments, the test set is never used for tuning.

\paragraph{Determinism and run metadata.}
Each run logs a \texttt{config.json} specifying (at minimum): random seed, dataset split paths, model hyperparameters, optimizer settings, number of steps/epochs, and evaluation settings (e.g., $K_{\mathrm{IS}}$ for IWAE). When feasible, we fix PyTorch and NumPy seeds and disable nondeterministic CuDNN behaviors. Reported results are either:
(i) single-run results with a fixed seed plus full artifacts, or
(ii) averages over multiple seeds with mean$\pm$std reporting.

\paragraph{Exported evaluation artifacts.}
All evaluation scripts export both raw per-sample metrics and summary tables:
\begin{itemize}
  \item \texttt{../results/metrics/}: CSV/JSON logs (per-step and per-epoch), per-image/per-sample scores, and summary tables.
  \item \texttt{../results/figures/}: plots used in the report (loss curves, KL curves, histograms, sample grids).
  \item \texttt{../results/samples/}: generated samples and traversals (VAE samples, GAN fixed-noise grids).
\end{itemize}
This ensures traceability: each plot/table in the report corresponds to a stored file.

% ---------------------------------------------------------
\subsection{VAE Metrics and Diagnostics}
We evaluate VAEs with complementary objectives: (i) predictive/likelihood-related evidence, (ii) representation and posterior diagnostics, and (iii) qualitative interpretability.

\paragraph{(V1) Loss decomposition and reporting conventions.}
We report the negative ELBO and its components:
\[
\mathcal{L}_{\mathrm{total}} = \mathcal{L}_{\mathrm{rec}} + \beta \cdot \mathcal{L}_{\mathrm{KL}},
\]
where $\mathcal{L}_{\mathrm{rec}}$ is the reconstruction term (BCE for Bernoulli pixels) and $\mathcal{L}_{\mathrm{KL}}$ is the KL divergence term. If a hierarchical latent model is used, we additionally report KL \emph{per level} (e.g., $\mathrm{KL}(z_1)$ and $\mathrm{KL}(z_2)$) and the total KL. In practice we log:
\begin{itemize}
  \item epoch/step, $\mathcal{L}_{\mathrm{rec}}$, $\mathcal{L}_{\mathrm{KL}}$ (and per-level KL if applicable), $\mathcal{L}_{\mathrm{total}}$,
  \item batch statistics (mean, median) and dataset-level aggregates on validation/test.
\end{itemize}
Interpretation: improvements in $\mathcal{L}_{\mathrm{total}}$ are only meaningful when both reconstruction and KL terms are analyzed, because a lower total loss can be achieved by collapsing KL to near-zero (posterior collapse) without learning useful latents.

\paragraph{(V2) Latent activity and active units.}
Latent utilization is measured using \emph{active units}. A standard criterion is:
\[
\text{AU}_j = \mathbb{I}\left(\mathrm{Var}_{x}\!\left[\mathbb{E}_{q_\phi(z\mid x)}(z_j)\right] > \tau_{\mathrm{AU}}\right),
\]
and the number of active units is $\sum_{j=1}^{d} \text{AU}_j$. We also report per-dimension KL contributions:
\[
\mathrm{KL}_j \approx \mathbb{E}_{x}\left[\mathrm{KL}\left(q_\phi(z_j\mid x)\,\|\,p(z_j)\right)\right],
\]
(or the corresponding quantity for a learned mixture prior).
Interpretation: a model with many inactive units (few AUs and near-zero $\mathrm{KL}_j$ for most dimensions) indicates collapse or over-regularization; a healthier latent representation exhibits non-trivial AUs and structured $\mathrm{KL}_j$.

\paragraph{(V3) Posterior--prior match (aggregate statistics).}
We compute aggregate posterior statistics over a large sample of datapoints:
\[
\mu_q = \mathbb{E}_{x}\left[\mathbb{E}_{q_\phi(z\mid x)}(z)\right],\qquad
\Sigma_q = \mathrm{Cov}_{x}\left(\mathbb{E}_{q_\phi(z\mid x)}(z)\right),
\]
and compare to:
(i) the standard normal prior ($\mu=0$, $\Sigma=I$), or
(ii) the learned prior (e.g., mixture implied by pseudo-inputs).
We also visualize $\Sigma_q$ in a PCA basis (eigen-spectrum and leading components) to assess whether the posterior occupies a lower-dimensional subspace.
Interpretation: strong mismatch can indicate either insufficient prior flexibility or an overly constrained posterior; a flexible prior (e.g., VampPrior) should reduce pathological mismatch while maintaining non-trivial latent structure.

\paragraph{(V4) Likelihood estimation via importance sampling (IWAE-style).}
Because ELBO is a lower bound, we estimate $\log p_\theta(x)$ using an importance-weighted estimator:
\[
\widehat{\log p_\theta(x)}=
\log\left(\frac{1}{K_{\mathrm{IS}}}\sum_{k=1}^{K_{\mathrm{IS}}}
\frac{p_\theta(x\mid z^{(k)})p(z^{(k)})}{q_\phi(z^{(k)}\mid x)}\right),
\quad z^{(k)}\sim q_\phi(z\mid x).
\]
For fairness, we use the same $K_{\mathrm{IS}}$ across models, and we report $K_{\mathrm{IS}}$ explicitly. We compute dataset-level averages and uncertainty estimates (e.g., mean$\pm$std across datapoints or seeds).
Interpretation: higher (less negative) values indicate better likelihood, but must be interpreted alongside collapse diagnostics, because some architectures can optimize likelihood while learning poor latent semantics.

\paragraph{(V5) Qualitative samples and latent traversals.}
We export:
\begin{itemize}
  \item unconditional samples: $z\sim p(z)$ then $x\sim p_\theta(x\mid z)$,
  \item latent traversals: vary one latent dimension while holding others fixed,
  \item PCA traversals: traverse along principal directions of inferred latent codes.
\end{itemize}
Interpretation: meaningful traversals suggest disentangled or at least interpretable factors; random/noisy traversals with low KL often correlate with collapse.

\paragraph{(V6) Pseudo-input visualization (VampPrior-specific).}
For VampPrior models, we visualize pseudo-inputs $\{u_k\}$ (after any range constraint such as $\tanh$) and optionally their decoded reconstructions. We interpret whether pseudo-inputs cover multiple modes (diverse digit styles/clothing prototypes).
Interpretation: diverse pseudo-inputs are consistent with a multimodal learned prior that can reduce over-regularization; highly similar pseudo-inputs can indicate poor utilization of the mixture capacity.

% ---------------------------------------------------------
\subsection{GAN Training Metrics (Conditional WGAN-GP)}
We evaluate GAN training using both \emph{objective-based} statistics and \emph{artifact-based} qualitative monitoring.

\paragraph{(G1) Critic and generator losses (decomposed).}
We log:
\begin{itemize}
  \item Wasserstein term: $\mathbb{E}[D(x_{\mathrm{real}},y)]-\mathbb{E}[D(x_{\mathrm{fake}},y)]$,
  \item critic total loss: Wasserstein term plus regularization,
  \item generator loss: $-\mathbb{E}[D(x_{\mathrm{fake}},y)]$.
\end{itemize}
Interpretation: the Wasserstein estimate should not diverge; excessively large magnitudes, instability, or oscillations often indicate training pathologies or learning-rate mismatch.

\paragraph{(G2) Regularization terms and constraint satisfaction.}
We explicitly plot:
\begin{itemize}
  \item gradient penalty magnitude and its moving average,
  \item drift penalty contribution,
  \item embedding regularization magnitude (if used).
\end{itemize}
Interpretation: gradient penalty should remain in a reasonable range; persistent extreme GP values may imply critic gradients are poorly conditioned or that Lipschitz constraints are not being satisfied.

\paragraph{(G3) Fixed-noise, fixed-label sample grids (checkpoint comparable).}
We export sample grids at fixed training steps (e.g., $\{0,1000,2000,3000,4000\}$ or assignment-specified checkpoints) using fixed noise vectors and a fixed label pattern covering all classes. This ensures that visual differences across grids reflect \emph{model changes} rather than sampling noise.
Interpretation: improvements should manifest as sharper class-appropriate structure and reduced mode collapse (diversity within each class).

\paragraph{(G4) Optional feature-space quality metric (FID-like proxy).}
If a consistent feature extractor is available (e.g., a fixed LeNet feature model), we compute a FID-like Fréchet distance between real and generated features. This is reported as a proxy quality measure and must be interpreted as \emph{feature-model dependent}.

% ---------------------------------------------------------
\subsection{Augmentation Evaluation (Classifier-Based)}
Augmentation is evaluated strictly on \emph{real} test data to measure generalization, not memorization.

\paragraph{(A1) Baseline vs augmented training protocol.}
We train the same classifier architecture under two regimes:
\begin{itemize}
  \item \textbf{Real-only baseline:} train on $\mathcal{D}_{\mathrm{CLS}}$ only.
  \item \textbf{Augmented:} train on $\mathcal{D}_{\mathrm{CLS}}$ plus $N_{\mathrm{syn}}$ synthetic samples per class generated by the trained conditional GAN.
\end{itemize}
All hyperparameters (optimizer, learning rate schedule, epochs, augmentations other than GAN) are held fixed to isolate the effect of GAN augmentation.

\paragraph{(A2) Performance endpoints and reporting.}
We report:
\begin{itemize}
  \item overall test accuracy on the untouched FashionMNIST test set,
  \item class-wise accuracy (or per-class error),
  \item optional confusion matrix to diagnose systematic class confusion,
  \item confidence intervals or mean$\pm$std across seeds (recommended in limited-data settings).
\end{itemize}
Interpretation: a scientifically meaningful improvement is one that is statistically consistent across seeds and not driven by a single class. Confusion matrices are particularly informative for detecting augmentation-induced boundary shifts.

\paragraph{(A3) Interpreting augmentation gains with sample quality diagnostics.}
Any observed accuracy improvement is interpreted jointly with:
\begin{itemize}
  \item conditional sample grids (class fidelity and diversity),
  \item critic diagnostics (stability and constraint satisfaction),
  \item class-wise performance (to detect whether some classes are harmed).
\end{itemize}
If augmentation improves accuracy but increases confusion in specific classes, we document this as a limitation and analyze whether the generator is producing ambiguous inter-class samples.

\paragraph{(A4) Failure modes and validity threats specific to augmentation.}
We explicitly document:
\begin{itemize}
  \item \textbf{Mode dropping:} generator produces limited intra-class diversity; can yield limited or misleading augmentation benefit.
  \item \textbf{Label noise via conditional confusion:} generator outputs samples that visually resemble a different class than the conditioning label, potentially harming classifier training.
  \item \textbf{Distribution shift amplification:} synthetic samples may overemphasize artifacts (checkerboarding, texture bias) that shift classifier decision boundaries.
\end{itemize}
Mitigation is evidence-driven: for example, reducing $N_{\mathrm{syn}}$, filtering low-quality samples, or improving conditional fidelity (architecture/regularization) is adopted only if it improves test performance and reduces class-wise harm.
